\documentclass[10pt,a4paper]{article}

\usepackage[margin=1.8cm]{geometry}
\usepackage[T1]{fontenc}
\usepackage[utf8]{inputenc}
\usepackage{charter}
\usepackage[english]{babel}
\usepackage{enumitem}
\usepackage{hyperref}
\usepackage{xcolor}
\usepackage{xurl}

\hypersetup{
colorlinks=true,
urlcolor=blue,
linkcolor=black
}

\setlength{\parindent}{0pt}
\setlength{\parskip}{0.45em}
\setlength{\emergencystretch}{2em}
\setlist[itemize]{leftmargin=1.4em,topsep=2pt,itemsep=2pt}

\newcommand{\relevance}[1]{\textbf{Industry relevance.} #1}

\begin{document}

\begin{center}
{\Large \textbf{Quantum and Computational Portfolio}} \\[3pt] 
{\large Antonio Michele Miti} \\[3pt]
\parbox{0.9\textwidth}{\centering
\textit{Selected projects and research outputs relevant to the
Quantum Data Scientist position at BIP AI Lab}}
\end{center}

\vspace{0.5em}

\section*{Profile}

I am a researcher with a background in \textbf{physics, mathematics, quantum algorithms, and scientific computing},
currently transitioning toward an industry-oriented role. My experience combines mathematical modelling with
\textbf{algorithm design, implementation, simulation, resource analysis, and quantitative comparison of alternative approaches}.

I am particularly interested in applying this background to \textbf{Quantum Computing Proofs of Concept, optimization,
simulation, and hybrid quantum--classical workflows}, where theoretical models must be translated into implementable
solutions and evaluated against practical constraints.

\section*{Core Quantum and Computational Experience}

\subsection*{1. Quantum arithmetic and resource optimization}

\textbf{A. Luongo, A. M. Miti, V. Narasimhachar, A. Sireesh,}\
\textit{Measurement-based uncomputation of quantum circuits for modular arithmetic},\
Proceedings of the 62nd ACM/IEEE Design Automation Conference (DAC '25), 2025.\
DOI: \href{https://doi.org/10.1109/DAC63849.2025.11132981}{10.1109/DAC63849.2025.11132981}.\\
arXiv: \href{https://arxiv.org/abs/2407.20167}{2407.20167}.

This work studies measurement-based uncomputation for modular-arithmetic circuits, focusing on reducing the resources
required by quantum arithmetic primitives.

\relevance{Direct experience with \textbf{quantum circuit design, uncomputation, modular arithmetic, resource estimation,
and implementation trade-offs}. The work required comparing alternative circuit constructions according to quantitative
costs.}

\subsection*{2. Hybrid CPU--GPU numerical implementation}

\textbf{Condensation Algorithm in CUDA}\
Public repository:
\href{https://github.com/MasterToninus/Computational-Physics-Exercise/tree/master/Condensation\%20Algorithm\%20in\%20CUDA}
{\texttt{Computational-Physics-Exercise / Condensation Algorithm in CUDA}}.

I developed a \textbf{hybrid CPU--GPU implementation of a determinant algorithm}, including a CUDA-based matrix library,
dedicated kernel tests, numerical validation, pivoting strategies, and comparison with CPU implementations.

\relevance{Hands-on experience with \textbf{C++, CUDA, GPU kernels, numerical algorithms, testing, and performance-oriented
implementation}. The project demonstrates the ability to translate a mathematical method into executable software,
validate its numerical behaviour, and compare alternative computational strategies.}

\subsection*{3. Monte Carlo simulation and algorithm comparison}

\textbf{Metropolis and cluster algorithms for the Ising model}\
Public repository:
\href{https://github.com/MasterToninus/Computational-Physics-Exercise/tree/master/Ising\%20simulation\%20Metropolis\%20vs\%20Cluster}
{\texttt{Computational-Physics-Exercise / Ising simulation}}.

This C++ project implements numerical experiments for the Ising model using \textbf{Metropolis and cluster-based
Monte Carlo algorithms}, together with supporting classes and numerical analysis.

\relevance{Experience with \textbf{simulation, stochastic algorithms, scientific programming, and benchmarking}.
In particular, it involved implementing alternative methods for the same physical problem and comparing their behaviour,
a useful basis for quantum--classical benchmarking and simulation-oriented projects.}

\section*{Mathematical Modelling and Problem Reduction}

\subsection*{4. Constraint handling and reduced models}

\textbf{A. M. Miti, L. Ryvkin,}\
\textit{Multisymplectic observable reduction using constraint triples},\
Differential Geometry and its Applications \textbf{100} (2025), 102272.\
DOI: \href{https://doi.org/10.1016/j.difgeo.2025.102272}{10.1016/j.difgeo.2025.102272}.\\
arXiv: \href{https://arxiv.org/abs/2506.00234}{2506.00234}.

The paper develops a systematic framework for reducing observables in constrained mathematical systems.

\relevance{Experience in \textbf{constraint analysis, symmetry reduction, and identification of effective degrees of
freedom}. These methods are transferable to optimization and simulation problems where reducing the dimensionality or
representation size can improve computational tractability.}

\subsection*{5. Structure-preserving model reduction}

\textbf{C. Blacker, A. M. Miti, L. Ryvkin,}\
\textit{Reduction of $L_\infty$-algebras of observables on multisymplectic manifolds},\
SIGMA \textbf{20} (2024), 061.\
DOI: \href{https://doi.org/10.3842/SIGMA.2024.061}{10.3842/SIGMA.2024.061}.\\
arXiv: \href{https://arxiv.org/abs/2206.03137}{2206.03137}.

This work studies how structured mathematical models behave under symmetry reduction.

\relevance{It demonstrates the ability to simplify a model while preserving the relations required for its interpretation.
In an industrial setting, this translates into identifying which variables, constraints, and dependencies must be retained
when moving from a theoretical formulation to an implementable model.}

\subsection*{6. Comparing alternative mathematical representations}

\textbf{A. M. Miti, M. Zambon,}\
\textit{Observables on multisymplectic manifolds and higher Courant algebroids},\
Communications in Contemporary Mathematics \textbf{27} (2025), no.~8, 2550006.\
DOI: \href{https://doi.org/10.1142/S0219199725500063}{10.1142/S0219199725500063}.\\
arXiv: \href{https://arxiv.org/abs/2209.05836}{2209.05836}.

The paper compares different mathematical descriptions of the same class of observables and establishes precise
relations between them.

\relevance{Experience in \textbf{representation choice and model comparison}. This is relevant when deciding how a
problem should be encoded for classical or quantum algorithms and when comparing alternative formulations according
to computational requirements.}

\section*{Additional Research Experience}

\textbf{D. Fiorenza, A. M. Miti,}\
\textit{$L_\infty$-morphisms between twisted Courant $r$-Lie algebras and untwisted Courant $(r+1)$-Lie algebroids},\
preprint, 2026.\
arXiv: \href{https://arxiv.org/abs/2602.14702}{2602.14702}.

This work develops explicit mappings between two higher-algebraic formulations and further demonstrates experience
with formal modelling and translation between mathematical representations.

\vspace{0.3em}

My broader peer-reviewed research also includes work on \textbf{symmetries, geometric quantization, conserved quantities,
and mathematical physics}. These projects document sustained experience in independent research, collaborative problem
solving, and the development of rigorous mathematical models.

\section*{Technical Skills and Industry Direction}

\textbf{Programming and computing:} Python, C++, CUDA, Linux, numerical methods, parallel computing, scientific software.

\textbf{Technical strengths:} algorithm design, numerical validation, simulation, resource analysis, benchmarking,
constraint handling, model reduction, and comparison of alternative computational approaches.

\textbf{Research and communication:} peer-reviewed publications, international collaborations, technical presentations,
seminars, teaching, and scientific documentation.

I am now intentionally moving toward an industrial environment where this combination of \textbf{quantum algorithms,
mathematical modelling, and scientific computing} can be applied to concrete technical problems. I am particularly
interested in projects where solutions are evaluated through \textbf{implementation, benchmarking, scalability,
resource constraints, and measurable performance} rather than only through theoretical analysis.

The Quantum Data Scientist role at BIP AI Lab is therefore a natural fit for my transition: it would allow me to apply
my current expertise to \textbf{Proof-of-Concept development, quantum--classical comparison, optimization, simulation,
and Data \& AI applications}, while further developing experience in industry-oriented project delivery.

\vfill

\noindent
\href{https://www.antoniomiti.it/}{www.antoniomiti.it}
\hfill
\href{https://orcid.org/0000-0002-8829-1943}{ORCID: 0000-0002-8829-1943}

\end{document}
